\section{Introduction to Automatic Image Analysis}

Automatic image analysis starts from the central difficulty of vision, then
moves through the main processing levels, the recurring goals of the field, and
the practical question of how to evaluate whether a vision system is actually
reliable.

\subsection{Lecture Aim and the Inverse Problem}

Automatic image analysis is the problem of turning measurements into
explanations. A digital image is easy to store: it is a grid of values. It is
much harder to understand: we want to infer objects, surfaces, motion, geometry,
lighting, and meaning from those values.

Computer vision is the computational study of how useful descriptions of the
visual world can be recovered from images and video. The description may be
semantic, such as ``there is a pedestrian in front of the car''; geometric, such
as ``this surface is two meters away''; or operational, such as ``reject this
part because it contains a defect''. This definition is deliberately broader
than object recognition. A vision system may label an image, reconstruct a
three-dimensional scene, align several images, estimate motion, improve an image
for a human observer, or supply measurements for a larger decision-making
system.

Computer vision is therefore not just image enhancement. It is inference under
uncertainty. The image is evidence; the scene that produced it is hidden.

\begin{figure}[ht]
  \centering
  \begin{tikzpicture}[
      node distance=2.2cm,
      box/.style={draw, rounded corners, align=center, minimum height=1cm,
        minimum width=2.9cm},
      arrow/.style={-{Latex[length=3mm]}, thick}
    ]
    \node[box, fill=blue!8] (scene) {3D scene\\objects, light, camera};
    \node[box, fill=green!8, right=of scene] (image) {2D image\\pixel grid};
    \node[box, fill=orange!10, right=of image] (meaning) {interpretation\\objects, depth, action};
    \draw[arrow] (scene) -- (image)
      node[midway, above=0.25cm, fill=white, inner sep=1pt]{image formation};
    \draw[arrow] (image) -- (meaning)
      node[midway, above=0.25cm, fill=white, inner sep=1pt]{inference};
    \draw[arrow, dashed] (image.south west) to[bend right=45]
      node[midway, below=0.35cm, fill=white, inner sep=1pt]{inverse problem}
      (scene.south east);
  \end{tikzpicture}
  \caption{Computer vision tries to invert the image formation process: pixels
  are observed, but the scene explanation is hidden.}
\end{figure}

\noindent
Human vision feels immediate, but this ease is misleading. A camera records a
two-dimensional projection of a three-dimensional world. During that projection,
many causes are mixed together: object shape, material, illumination, viewpoint,
occlusion, sensor noise, and motion. Different scenes can produce very similar
images, so there is usually no unique answer that can be read directly from the
pixels.

This makes vision an inverse problem. The forward direction is comparatively
clear: if we know the scene, the light, the camera, and the materials, we can
model how an image is formed. The inverse direction asks us to go backward from
image measurements to a plausible explanation of the world. That backward step
requires assumptions, models, data, and algorithms.

A compact way to write the forward model is
\begin{equation}
  I(x,y) = f\!\left(\Pi(X,Y,Z), L, M, C\right) + \varepsilon ,
\end{equation}
\noindent
where $I(x,y)$ is the measured pixel value, $\Pi$ is the camera projection from
3D to 2D, $L$ describes illumination, $M$ material and surface properties, $C$
camera response, and $\varepsilon$ sensor noise. Vision asks for some hidden
scene state $S$ from the image $I$:
\begin{equation}
  \hat{S} = \arg\max_S p(S \mid I)
  = \arg\max_S p(I \mid S)\,p(S).
\end{equation}
\noindent
The likelihood $p(I \mid S)$ says which images a scene would probably produce.
The prior $p(S)$ encodes what kinds of scenes are plausible before seeing the
image.

\begin{figure}[ht]
  \centering
  \begin{tikzpicture}[
      node distance=0.9cm,
      scene/.style={draw, trapezium, trapezium left angle=70,
        trapezium right angle=110, minimum width=2.4cm, minimum height=1cm,
        align=center, fill=blue!8},
      img/.style={draw, minimum width=2.4cm, minimum height=1cm,
        align=center, fill=green!8},
      arrow/.style={-{Latex[length=3mm]}, thick}
    ]
    \node[scene] (s1) {Scene A\\small nearby object};
    \node[scene, below=of s1] (s2) {Scene B\\large distant object};
    \node[img, right=2.4cm of s1, yshift=-0.95cm] (img) {similar\\image size};
    \draw[arrow] (s1.east) -- (img.west);
    \draw[arrow] (s2.east) -- (img.west);
  \end{tikzpicture}
  \caption{Ambiguity is unavoidable: different scenes can lead to similar image
  measurements. Priors and additional evidence help choose an explanation.}
\end{figure}

The gap between the pixel array and the human interpretation is the essential
problem. A person immediately infers objects, depth, material, intent, and
context, while the measured image itself only contains sampled intensities or
colors.

\subsection{From Pixels to Interpretation}

Automatic image analysis can be organized as a progression from low-level to
high-level tasks.

\begin{itemize}
  \item Low-level processing improves or transforms image measurements, for
        example acquisition, sampling, denoising, filtering, contrast changes,
        and edge extraction. Its typical mapping is image to image, or scene to
        image when the acquisition process is included.
  \item Mid-level analysis groups measurements into more meaningful structures,
        such as regions, contours, features, segments, motion fields, or depth
        estimates. It is often the bridge from image measurements to feature
        descriptions.
  \item High-level interpretation assigns semantic meaning, such as recognizing
        objects, estimating scene layout, detecting abnormalities, or tracking
        actions. Its typical mapping is from features or image evidence to
        object models, object detections, categories, or decisions.
\end{itemize}

\noindent
These levels are useful for teaching, but real systems often mix them. A modern
recognition model may learn low-level filters and high-level categories inside
one network. A classical pipeline may use carefully designed features followed
by statistical inference. In both cases, the goal is the same: reduce the
ambiguity of image data until a useful decision or description becomes possible.

\begin{table}[ht]
  \centering
  \begin{tabular}{@{}p{2.0cm}p{3.0cm}p{4.2cm}p{4.0cm}@{}}
    \toprule
    Level & Typical mapping & Typical question & Examples \\
    \midrule
    Low-level & scene/image $\rightarrow$ image &
      What local measurements are reliable? &
      denoising, gradients, edges \\
    Mid-level & image $\rightarrow$ features &
      Which measurements belong together? &
      regions, contours, segments, motion \\
    High-level & features/image $\rightarrow$ objects &
      What does the image mean? &
      objects, scene layout, actions \\
    \bottomrule
  \end{tabular}
  \caption{A teaching taxonomy for image-analysis tasks.}
\end{table}

\noindent
The same hierarchy can also be read as a contrast between bottom-up and
top-down processing. A bottom-up system starts with the image data and builds
larger hypotheses from local evidence: pixels become edges, edges become
contours, contours become regions, and regions become object candidates. This
is data-driven and attractive because the image itself constrains the result.
Its weakness is that local evidence is often ambiguous or missing.

Top-down processing starts from a model, context, task goal, or prior
expectation and asks how the image should be interpreted under that hypothesis.
This is knowledge-driven: a rough object model can guide the search for parts,
a scene context can make one interpretation more plausible than another, and a
classifier can project learned category knowledge onto the image. Top-down
processing is powerful because it uses what is already known, but it can also
produce confident misinterpretations when the prior is wrong. Practical systems
therefore often use a mixed mode: bottom-up evidence proposes or constrains
hypotheses, while top-down models select, refine, or reject them.

\begin{figure}[ht]
  \centering
  \begin{tikzpicture}[
      node distance=0.85cm,
      box/.style={draw, rounded corners, align=center, minimum width=2.75cm,
        minimum height=0.9cm},
      arrow/.style={-{Latex[length=3mm]}, thick}
    ]
    \node[box, fill=green!8] (pixels) {pixels};
    \node[box, fill=green!8, right=of pixels] (features) {features};
    \node[box, fill=green!8, right=of features] (parts) {parts or regions};
    \node[box, fill=orange!12, right=of parts] (object) {object hypothesis};
    \node[box, fill=blue!8, above=1.0cm of parts] (model) {model, context, task};

    \draw[arrow] (pixels) -- (features);
    \draw[arrow] (features) -- (parts);
    \draw[arrow] (parts) -- (object);
    \draw[arrow] (model) -- (features)
      node[midway, left=0.12cm, fill=white, inner sep=1pt, font=\small]
      {top-down};
    \draw[arrow] (model) -- (object);
    \node[align=center, below=0.2cm of features, font=\small]
      {bottom-up evidence};
  \end{tikzpicture}
  \caption{Bottom-up processing builds interpretations from image evidence.
  Top-down processing uses models and context to guide or correct the
  interpretation.}
\end{figure}

\noindent
One common low-level operation is convolution:
\begin{equation}
  (I * k)(x,y) = \sum_{u=-r}^{r}\sum_{v=-r}^{r} I(x-u,y-v)\,k(u,v).
\end{equation}
\noindent
The same mathematical form can blur, sharpen, estimate derivatives, or detect
local structure depending on the kernel $k$.

Computer vision can be organized around three recurring goals: recognition,
reconstruction, and reorganization.

\paragraph{Recognition.}
Recognition asks what is present and where it is. Examples include
classification, detection, segmentation, reading text, identifying faces, or
finding defects in industrial inspection.

\paragraph{Reconstruction.}
Reconstruction asks what spatial or physical structure produced the image.
Examples include estimating depth, camera motion, three-dimensional shape, or
the geometry of a scene from one or more images.

\paragraph{Reorganization.}
Reorganization changes the image representation so that later analysis becomes
easier or the image becomes more useful to a person. Examples include
registration, stitching, stabilization, denoising, super-resolution, and
computational photography.

These goals are connected. Recognition benefits from good representations;
reconstruction needs reliable correspondences; reorganization often depends on
understanding motion, geometry, or image formation.

\begin{table}[ht]
  \centering
  \begin{tabular}{@{}p{3.3cm}p{4.8cm}p{5.2cm}@{}}
    \toprule
    Area & Main transformation & Typical examples \\
    \midrule
    Digital image processing & image $\rightarrow$ image & enhancement,
      restoration, filtering, compression \\
    Photogrammetry and geometry & image(s) $\rightarrow$ 3D model & camera
      calibration, depth, surface reconstruction, mapping \\
    Automatic image analysis & image $\rightarrow$ object model or detection &
      feature extraction, segmentation, object categorization, recognition \\
    \bottomrule
  \end{tabular}
  \caption{Automatic image analysis is closest to the recognition branch of
  computer vision: it studies how image evidence becomes object-level
  descriptions, while still using processing and geometric ideas when they
  support that goal.}
\end{table}

\begin{figure}[ht]
  \centering
  \begin{tikzpicture}[
      goal/.style={circle, draw, minimum size=2.8cm, align=center},
      note/.style={align=center},
      arrow/.style={-{Latex[length=3mm]}, thick}
    ]
    \node[goal, fill=red!8] (rec) {Recognition\\what and where?};
    \node[goal, fill=blue!8, right=2.8cm of rec] (recon)
      {Reconstruction\\what 3D structure?};
    \node[goal, fill=green!8, below=1.8cm of $(rec)!0.5!(recon)$] (reorg)
      {Reorganization\\better representation};
    \draw[arrow] (rec) -- (recon);
    \draw[arrow] (recon) -- (reorg);
    \draw[arrow] (reorg) -- (rec);
  \end{tikzpicture}
  \caption{Recognition, reconstruction, and reorganization are not isolated
  topics; each goal supplies information that can support the others.}
\end{figure}

\noindent
Another useful taxonomy separates the vocabulary of the field by the kind of
question being asked. Some terms describe the \emph{input} and its formation,
some describe \emph{representations} computed from the image, some describe
\emph{tasks}, and some describe the \emph{knowledge} used to resolve ambiguity.
Keeping these groups separate helps avoid a common beginner confusion: a feature,
a model, a task, and an application are not the same kind of object.

\begin{figure}[ht]
  \centering
  \resizebox{\textwidth}{!}{%
  \begin{tikzpicture}[
      node distance=0.85cm and 0.95cm,
      root/.style={draw, rounded corners, align=center, minimum width=3.4cm,
        minimum height=0.9cm, fill=gray!10},
      group/.style={draw, rounded corners, align=center, minimum width=3.5cm,
        minimum height=0.9cm},
      item/.style={draw, align=center, minimum width=3.1cm, minimum height=0.72cm,
        font=\small},
      arrow/.style={-{Latex[length=2.4mm]}, thick}
    ]
    \node[root] (cv) {computer vision\\images to explanations};

    \node[group, fill=blue!8, below left=1.0cm and 5.5cm of cv] (input)
      {image formation};
    \node[group, fill=green!8, below left=1.0cm and 1.7cm of cv] (repr)
      {representations};
    \node[group, fill=orange!12, below right=1.0cm and 1.7cm of cv] (tasks)
      {tasks};
    \node[group, fill=red!8, below right=1.0cm and 5.5cm of cv] (knowledge)
      {knowledge};

    \node[item, below=of input] (inputa) {scene, light, camera};
    \node[item, below=of inputa] (inputb) {projection, sampling, noise};

    \node[item, below=of repr] (repra) {pixels, color, gradients};
    \node[item, below=of repra] (reprb) {features, regions, motion};
    \node[item, below=of reprb] (reprc) {depth, pose, 3D shape};

    \node[item, below=of tasks] (taska) {recognition};
    \node[item, below=of taska] (taskb) {reconstruction};
    \node[item, below=of taskb] (taskc) {reorganization};

    \node[item, below=of knowledge] (knowa) {geometry and physics};
    \node[item, below=of knowa] (knowb) {probability and priors};
    \node[item, below=of knowb] (knowc) {learned models and data};

    \draw[arrow] (cv) -- (input);
    \draw[arrow] (cv) -- (repr);
    \draw[arrow] (cv) -- (tasks);
    \draw[arrow] (cv) -- (knowledge);
    \draw[arrow] (input) -- (inputa);
    \draw[arrow] (inputa) -- (inputb);
    \draw[arrow] (repr) -- (repra);
    \draw[arrow] (repra) -- (reprb);
    \draw[arrow] (reprb) -- (reprc);
    \draw[arrow] (tasks) -- (taska);
    \draw[arrow] (taska) -- (taskb);
    \draw[arrow] (taskb) -- (taskc);
    \draw[arrow] (knowledge) -- (knowa);
    \draw[arrow] (knowa) -- (knowb);
    \draw[arrow] (knowb) -- (knowc);
  \end{tikzpicture}%
  }
  \caption{A compact taxonomy of introductory computer-vision terms. The field
  connects image formation, intermediate representations, task goals, and prior
  knowledge.}
\end{figure}

\subsection{Models, Data, and Evaluation}

Because images are ambiguous, vision systems need prior knowledge. Sometimes
that knowledge is explicit: camera geometry, physical image formation, smooth
surfaces, rigid motion, or probabilistic noise models. Sometimes it is learned
from data: neural networks discover features and decision rules from many
examples.

Classical computer vision often emphasizes explicit models and engineered
features. Modern deep learning often emphasizes learned representations and
large datasets. The practical lesson is not that one replaces the other
completely. Robust systems usually combine several kinds of knowledge:
scientific models for image formation, statistical models for uncertainty,
engineering choices for reliability, and data-driven learning for complex
patterns that are hard to specify by hand.

\begin{figure}[ht]
  \centering
  \resizebox{\textwidth}{!}{%
  \begin{tikzpicture}[
      milestone/.style={draw, rounded corners, align=center, minimum width=3cm,
        minimum height=1.05cm, font=\small},
      arrow/.style={-{Latex[length=3mm]}, thick}
    ]
    \node[milestone, fill=gray!10] (m1960) {1960s--1970s\\blocks world\\edges, line labeling};
    \node[milestone, fill=blue!8, right=0.7cm of m1960] (m1980)
      {1980s\\quantitative vision\\pyramids, shape-from-X};
    \node[milestone, fill=green!8, right=0.7cm of m1980] (m1990)
      {1990s\\multiple views\\stereo, motion, segmentation};
    \node[milestone, fill=orange!12, right=0.7cm of m1990] (m2000)
      {2000s\\features and graphics\\recognition, stitching};
    \node[milestone, fill=red!8, right=0.7cm of m2000] (m2010)
      {2010s--today\\deep learning\\large data, sensors};
    \draw[arrow] (m1960) -- (m1980);
    \draw[arrow] (m1980) -- (m1990);
    \draw[arrow] (m1990) -- (m2000);
    \draw[arrow] (m2000) -- (m2010);
  \end{tikzpicture}%
  }
  \caption{A brief historical timeline. The emphasis moved from hand-built
  geometric reasoning, through quantitative reconstruction and engineered
  features, toward large-scale learned representations and integrated sensing.}
\end{figure}

\noindent
The timeline is not a story of one idea simply replacing another. Edge
operators, camera geometry, variational estimation, probabilistic models,
feature matching, and neural networks all remain useful. What changes over time
is the balance between explicit modeling, available computation, available
training data, and the kinds of applications that can be deployed reliably.

\begin{table}[ht]
  \centering
  \begin{tabular}{@{}p{3.1cm}p{5cm}p{5cm}@{}}
    \toprule
    Source of knowledge & What it contributes & Risk if missing \\
    \midrule
    Physical model & Camera geometry, illumination, projection & Results may
      violate how images are formed \\
    Statistical model & Noise, uncertainty, priors, inference & The system may
      overtrust ambiguous measurements \\
    Engineering design & Robustness, speed, failure handling & The method may
      work only in demonstrations \\
    Data-driven learning & Features and decision rules from examples & Complex
      visual variation may be underspecified \\
    \bottomrule
  \end{tabular}
  \caption{A practical vision system usually combines several kinds of
  knowledge.}
\end{table}

\noindent
Computer vision is now a basic technology in many domains. It is more useful to
name the application categories than to memorize a long list of products:
consumer imaging and media, medicine and life sciences, industry and logistics,
robotics and autonomous systems, remote sensing and mapping, document analysis,
security and biometrics, and scientific measurement. The tasks differ, but the
same conceptual questions recur: what image evidence is available, what
assumptions make the interpretation possible, and what failure cases matter in
the application context?

\begin{table}[ht]
  \centering
  \begin{tabular}{@{}p{3.6cm}p{4.8cm}p{5.1cm}@{}}
    \toprule
    Application category & Typical vision output & Main reliability concern \\
    \midrule
    Consumer and media & enhanced images, alignment, tracking & visual quality
      and robustness across devices \\
    Medicine and life sciences & measurements, segmentations, suspected findings
      & uncertainty, bias, and clinical validation \\
    Industry and logistics & defects, codes, counts, object poses & speed,
      repeatability, and rare failures \\
    Robotics and autonomy & localization, obstacles, maps, actions & safety
      under changing environments \\
    Remote sensing and mapping & land cover, 3D structure, change detection &
      scale, resolution, and registration accuracy \\
    Documents and biometrics & text, identity cues, authenticity signals &
      privacy, fairness, and adversarial misuse \\
    \bottomrule
  \end{tabular}
  \caption{Categorical application areas are easier to study than isolated
  product examples because they expose recurring reliability questions.}
\end{table}

These applications also show why evaluation matters. A result that looks
plausible on one example is not enough. Vision algorithms must be tested against
noise, changing illumination, clutter, unusual viewpoints, occlusion, and data
that differ from the training or development examples.

\begin{figure}[ht]
  \centering
  \begin{tikzpicture}[
      evalstep/.style={draw, rounded corners, align=center, minimum width=3.2cm,
        minimum height=1cm},
      arrow/.style={-{Latex[length=3mm]}, thick}
    ]
    \node[evalstep, fill=green!8] (clean) {clean synthetic data\\known truth};
    \node[evalstep, fill=yellow!12, right=of clean] (noise) {controlled stress\\noise, blur, light};
    \node[evalstep, fill=red!8, right=of noise] (real) {real images\\clutter, bias, surprises};
    \draw[arrow] (clean) -- (noise);
    \draw[arrow] (noise) -- (real);
  \end{tikzpicture}
  \caption{Evaluation should move from controlled correctness toward real-world
  robustness.}
\end{figure}

\begin{figure}[ht]
  \centering
  \begin{tikzpicture}[
      node distance=0.7cm,
      block/.style={draw, rounded corners, align=center, minimum width=2.9cm,
        minimum height=0.95cm},
      small/.style={draw, rounded corners, align=center, minimum width=2.35cm,
        minimum height=0.85cm},
      arrow/.style={-{Latex[length=3mm]}, thick}
    ]
    \node[block, fill=blue!8] (pixels) {image measurements\\pixels, gradients, features};
    \node[block, fill=orange!10, right=of pixels] (models)
      {prior knowledge\\physics, statistics, data};
    \node[block, fill=green!8, right=of models] (output)
      {useful output\\label, geometry, improved image};

    \node[small, fill=red!8, below=1.1cm of pixels] (rec)
      {recognition\\what is where?};
    \node[small, fill=blue!8, below=1.1cm of models] (recon)
      {reconstruction\\what 3D scene?};
    \node[small, fill=green!8, below=1.1cm of output] (reorg)
      {reorganization\\better representation};

    \draw[arrow] (pixels) -- (models);
    \draw[arrow] (models) -- (output);
    \draw[arrow] (pixels) -- (rec);
    \draw[arrow] (models) -- (recon);
    \draw[arrow] (output) -- (reorg);
  \end{tikzpicture}
  \caption{A compact summary of the introductory framing: measurements become
  useful outputs only after they are combined with assumptions, models, and
  data.}
\end{figure}

\subsection{Topic Roadmap}

Automatic image analysis is narrower than the whole field of computer vision.
It uses image processing tools and needs some geometric and probabilistic
reasoning, but the main line is automatic interpretation: how images are
transformed into features, segments, object models, detections, categories, and
evaluated decisions.

The foundations are image representation, filtering, features, transformations,
and the basic language of image formation. These topics supply the bottom-up
evidence used later. Segmentation, clustering, and feature extraction provide
mid-level descriptions that make object-level reasoning possible. Learning-based
methods, especially neural networks, show how useful representations and
classifiers can be estimated from data. Object models, probabilistic estimation,
and higher-level recognition tasks connect these pieces; reconstruction appears
where it helps explain image formation, geometry, or the limits of recognition.

Three questions recur throughout automatic image analysis:

\begin{enumerate}
  \item What information is actually present in the image measurements?
  \item What assumptions are needed to infer the hidden scene or decision?
  \item How can we test whether the algorithm works beyond the example in front
        of us?
\end{enumerate}

\noindent
Automatic image analysis starts with pixels but does not end there. Its real
task is to connect images to explanations: what is in the scene, how it is
arranged, how it changes, and what action or decision should follow.
